% =====================================================================
%  c13：原印刷页 107–113（PDF p.108–114）
%  内容：§4.1.2 全节（含七条编号公式 (4.32)–(4.47)）；原 p.109 的练习与
%        Remarks 列表（(4.48)–(4.50)）；§4.2 及其前段（(4.51)–(4.73)：
%        三种推导途径中的第一种、第二种、第三种）。
%  本片段以 §4.1.2 的标题起头；§4.2 按编号点数（两点）判为 \subsection，
%  故全片段无 \section 级标题，不插入 keytakeaways。本片段无插图。
%  编号公式 (4.32)–(4.73) 连续，无缺号。
% =====================================================================

\subsubsection{二维情形与共形点} % 原文: 4.1.2 The Case Of Two Dimensions And The Conformal Point

在高于二维的情形，WZ 项会给运动方程带来高阶导数的修正。我们必须把它理解为一个有效场论（effective field theory）。在二维，情况不同，而且极其优美：对耦合常数 $f$ 的特殊取值，运动方程具有无穷维对称性。

我们强烈建议读者参阅那两篇解决了量子理论的优美论文：

\begin{enumerate}
  \item E.~Witten, ``Nonabelian bosonization in two dimensions,'' Commun.\ Math.\ Phys.\ 92:455--472, 1984.
  \item V.~G.~Knizhnik and A.~B.~Zamolodchikov, ``Current Algebra and Wess-Zumino Model in Two-Dimensions,'' Nucl.\ Phys.\ B247:83--103, 1984.
\end{enumerate}

当 $f^2=12\pi k c_1$ 时，作用量成为 $S=kS^+_{WZW}$，其中
\begin{equation}
  S^+_{WZW} = 6\pi c_1\int\operatorname{Tr}\left(g^{-1}\partial_+ g\,g^{-1}\partial_- g\right)dx^+\wedge dx^- + 2\pi c_1\int_B\operatorname{Tr}\left(g^{-1}dg\right)^3
  \tag{4.32}
  \label{eq:4.32}
\end{equation}
% [Tutor Note: (4.32) 的归一化是全片段所有系数的源头，最省事的核对办法是把它与 p.110 的 (4.50) 对齐：对 $G=SU(N)$ 取 $N$ 表示中的迹时 $c_1=1/(24\pi^2)$，故 $S=kS^+_{WZW}$ 中 WZ 项的系数为 $2\pi c_1k=\frac{k}{12\pi}$，即 (4.50) 指数上的量；而 (4.32) 的第一项给出 $\frac{k}{4\pi}\int\operatorname{Tr}(g^{-1}\partial_+ g\,g^{-1}\partial_- g)$，正是熟知的 WZW 作用量归一化。若把 $c_1$ 记成 $1/(24\pi)$ 或把 $12\pi kc_1$ 写成 $24\pi kc_1$，(4.35)、(4.36) 与 (4.53) 里的 $\partial_\pm$ 前因子会整体差一个 2，而 (4.59) 给出的能级 $k$ 也随之错位。]

而若 $f^2=-12\pi k c_1$，则作用量为 $S=-kS^-_{WZW}$，其中
\begin{equation}
  S^-_{WZW} = 6\pi c_1\int\operatorname{Tr}\left(g^{-1}\partial_+ g\,g^{-1}\partial_- g\right)dx^+\wedge dx^- - 2\pi c_1\int_B\operatorname{Tr}\left(g^{-1}dg\right)^3
  \tag{4.33}
  \label{eq:4.33}
\end{equation}

这些作用量满足一些漂亮的性质：见下面关于 Polyakov-Wiegmann（PW）公式的练习。这些恒等式有着重要的物理含义。

在二维，运动方程为：
\begin{equation}
  -\frac{f^2}{2}d\left(\ast g^{-1}dg\right) + 6\pi k c_1\left(g^{-1}dg\right)^2 = 0
  \tag{4.34}
  \label{eq:4.34}
\end{equation}

对第二项使用 MC 方程，我们可以把它写成：
\begin{equation}
  d\left(\frac{f^2}{2}\ast g^{-1}dg + 6\pi k c_1 g^{-1}dg\right) = 0
  \tag{4.35}
  \label{eq:4.35}
\end{equation}
\begin{equation}
  d\left(\ast g^{-1}dg + \frac{12\pi k c_1}{f^2}g^{-1}dg\right) = 0
  \tag{4.36}
  \label{eq:4.36}
\end{equation}
% [Tutor Note: 原文「对第二项使用 MC 方程」一语带过了唯一的一步代数。Maurer-Cartan 方程即 $d(g^{-1}dg)=-(g^{-1}dg)^2$，故 $(g^{-1}dg)^2=-d(g^{-1}dg)$；代入 (4.34) 得 $-\frac{f^2}{2}d(\ast g^{-1}dg)-6\pi kc_1\,d(g^{-1}dg)=0$，整体乘 $-2/f^2$ 并把 $d$ 提到括号外，正是 (4.35)；(4.36) 只是把 (4.35) 除以 $f^2/2$。注意 (4.34) 的第二项是 $(g^{-1}dg)^2$ 而不是 $d(g^{-1}dg)$，符号一旦弄反，得到的将是 $d(\ast g^{-1}dg - \frac{12\pi kc_1}{f^2}g^{-1}dg)=0$，后面 (4.40) 与 (4.41) 的手征性也会整体互换。]

现在考虑度规为 $ds^2=-dx^+dx^-=-dt^2+dx^2$ 的 Minkowski 空间。我们可以写
\begin{equation}
  g^{-1}dg = \left(g^{-1}\partial_+ g\right)dx^+ + \left(g^{-1}\partial_- g\right)dx^-
  \tag{4.37}
  \label{eq:4.37}
\end{equation}

于是，利用定向 $dt\wedge dx$，自对偶形式为
\begin{equation}
  \begin{aligned}
    \ast(dt+dx) &= \ast dx^+ = -dx^+\\
    \ast(dt-dx) &= \ast dx^- = dx^-
  \end{aligned}
  \tag{4.38}
  \label{eq:4.38}
\end{equation}
% [Tutor Note: (4.38) 隐含了 Lorentz 号差下二维 Hodge 星的全部约定：由 $\ast dt=-dx$、$\ast dx=-dt$ 可得 $\ast(dt+dx)=-dx-dt=-dx^+$ 与 $\ast(dt-dx)=-dx+dt=dx^-$，且 $\ast^2=-1$（对 1-形式），与欧氏号差下的 $\ast^2=+1$ 相反。若误用欧氏约定 $\ast dt=dx$、$\ast dx=dt$，则 (4.38) 两行都差一个符号，(4.39) 的自对偶组合随之变成 $\ast g^{-1}dg=(g^{-1}\partial_+ g)dx^+-(g^{-1}\partial_- g)dx^-$，代入 (4.36) 后得到的将是 $\partial_-(g^{-1}\partial_+ g)=0$ 而非 (4.40)——也就是把左行与右行对调。]

因此
\begin{equation}
  \ast g^{-1}dg = -\left(g^{-1}\partial_+ g\right)dx^+ + \left(g^{-1}\partial_- g\right)dx^-
  \tag{4.39}
  \label{eq:4.39}
\end{equation}

由此可见
\begin{itemize}
  \item 若 $f^2=12\pi k c_1$，则运动方程为
  \begin{equation}
    \partial_+\left(g^{-1}\partial_- g\right) = 0
    \tag{4.40}
    \label{eq:4.40}
  \end{equation}
  \item 若 $f^2=-12\pi k c_1$，则运动方程为
  \begin{equation}
    \partial_-\left(g^{-1}\partial_+ g\right) = 0
    \tag{4.41}
    \label{eq:4.41}
  \end{equation}
\end{itemize}

我们现在看到这个理论的一个非凡性质。为确定起见，让我们集中看 $f^2=12\pi k c_1$ 的情形。我们可以定义无穷维群
\begin{equation}
  \mathcal{G} := Map(\mathbb{R}, G)
  \tag{4.42}
  \label{eq:4.42}
\end{equation}
于是存在 $\mathcal{G}_L\times\mathcal{G}_R$ 在运动方程解集上的左作用！

事实上，若 $g(x^+,x^-)$ 解出 (4.40)，则下式亦然
\begin{equation}
  h_l(x^+)g(x^+,x^-)h_r(x^-)^{-1}
  \tag{4.43}
  \label{eq:4.43}
\end{equation}

一个与之密切相关的论点是：运动方程蕴含存在两条\emph{分别}守恒的流。定义
\begin{equation}
  J_- := g^{-1}\partial_- g
  \tag{4.44}
  \label{eq:4.44}
\end{equation}
\begin{equation}
  J_+ := \partial_+ g\,g^{-1}
  \tag{4.45}
  \label{eq:4.45}
\end{equation}
于是 (4.40) 连同恒等式
\begin{equation}
  \partial_+\left(g^{-1}\partial_- g\right) = g^{-1}\partial_-\left(\partial_+ g\,g^{-1}\right)g
  \tag{4.46}
  \label{eq:4.46}
\end{equation}
表明 $\partial_+J_-=0$ 与 $\partial_-J_+=0$ 同时成立。

所以当 $f^2=12\pi k c_1$ 时有两条分别守恒的流。

让我们借此求出运动方程的通解。我们必须选取某种边界条件。

若在 $\mathbb{R}^{1,1}$ 上工作，则当 $x^+\to-\infty$ 时 $J_+(x^+)\to 0$ 足够快，使得 $\int^{x^+}_{-\infty}J_+$ 收敛，$J_-$ 亦然。于是我们可以把 $J_\pm$ 视为给定，并积分方程 (4.44)(4.45) 而得到\footnote{我们对路径排序的约定是把较晚的时刻放在左边。}
\begin{equation}
  g(x^+,x^-) = \left(P\exp\int^{x^+}_{-\infty}J_+(s)ds\right)\cdot g_0\cdot\left(P\exp -\int^{x^-}_{-\infty}J_-(s)ds\right)^{-1}
  \tag{4.47}
  \label{eq:4.47}
\end{equation}
% [Tutor Note: (4.47) 有三处极易写错：其一，路径排序的约定见脚注 33（较晚时刻在左），故两个 $P\exp$ 的因子次序不可交换；其二，右因子必须整体取 $-1$ 次幂：由于 $J_-=g^{-1}\partial_- g$ 是右不变形式，解出 $g$ 时得到的是 $\partial_- h=hJ_-$（$J_-$ 在右），其解是反路径排序的 $\bar P\exp(\int J_-)$，而 $\left(P\exp(-\int J_-)\right)^{-1}$ 恰好等于它；漏掉 $-1$ 次幂就既还原不出 $g=1$ 的平凡解，也违反 (4.44)。其三，左因子含 $J_+$（$x^+$ 方向）、右因子含 $J_-$（$x^-$ 方向），与 (4.44)(4.45) 中 $J_\pm$ 的定义方向一致，切勿对调。$g_0$ 是 $J_\pm\to 0$ 处残留的常值群元。]

\emph{换句话说，在 $\mathbb{R}^{1,1}$ 上，运动方程的通解是群 $\mathcal{G}_L\times\mathcal{G}_R$ 对 $g=1$ 的一个变换。}

% 原 p.108 页末另有一行星号分隔线（作者分隔草稿段落所用），此处不转写。

{\small\itshape ［原文标记：需要讨论欧氏延拓——全纯流与反全纯流］}
% [原文待办标记，原文 p.109 页首，原文全大写并以下方一行星号包裹；此处照原文保留为小号斜体]

\medskip
\noindent\textbf{练习}\quad 写出 $f^2=-12\pi k c_1$ 情形的相应方程。

\medskip
\noindent\textbf{练习}（\emph{PW 公式的一个推论}）\quad 用 PW 公式证明
\begin{equation}
  S^+_{WZW}(g_1g_2) = S^+_{WZW}(g_1) + S^+_{WZW}(g_2) + 12\pi c_1\int\operatorname{Tr}\left(\partial_+ g_2\,g_2^{-1}g_1^{-1}\partial_- g_1\right)
  \tag{4.48}
  \label{eq:4.48}
\end{equation}
\begin{equation}
  S^-_{WZW}(g_1g_2) = S^-_{WZW}(g_1) + S^-_{WZW}(g_2) + 12\pi c_1\int\operatorname{Tr}\left(\partial_- g_2\,g_2^{-1}g_1^{-1}\partial_+ g_1\right)
  \tag{4.49}
  \label{eq:4.49}
\end{equation}

{\small\itshape ［原文标记：需要与上文把 $A_{\bar z}$ 规范到零处关于高斯定律的讨论作比较］}
% [原文待办标记，原文 p.109，上下各有一行星号分隔线]

\medskip
\noindent\textbf{练习}\quad 对二维 $SL(2,\mathbb{R})$ WZW 模型，$k$ 是否量子化？

\medskip
\noindent\textbf{Remarks（评注）：}

\begin{itemize}
  \item 注意「$J_+$ 是左对称性」这句话有\emph{两重}含义。它描述左行波，因为它是 $x^+=t+x$ 的函数，故随时间向前推进，波向负 $x$ 方向运动。它同时还对应于把 $g(x^+,x^-)$ 从左边乘上某个群元。
  \item \emph{推广到欧氏号差。} 理论可以延拓到欧氏号差 $x^+\to z$、$x^-\to\bar z$。此时对称性是全纯/反全纯流代数。由于共形不变性，理论可以在任意 Riemann 面上表述。
  \item 我们已经证明耦合常数 $k$ 是拓扑量子化的。于是问题来了：「$k$ 的最小值是多少？」这取决于表示以及归一化常数 $c_n$。对二维情形，我们有 $SU(N)$ 的 WZ 项：
  \begin{equation}
    \exp\left[ik\frac{1}{12\pi}\int_{B_3}\operatorname{Tr}_N\left(g^{-1}dg\right)^3\right]
    \tag{4.50}
    \label{eq:4.50}
  \end{equation}
  其中 $k\in\mathbb{Z}$，并且所有整数都能出现。因此，对 $G=SU(N)$ 取 $N$ 表示中的迹，我们有 $c_1=\frac{1}{24\pi^2}$。
  \item \emph{谱流（spectral flow）与运动方程的有趣解。} $SL(2,\mathbb{R})$ 中的长弦（long string）等。
\end{itemize}

\subsection{量子化：Virasoro 与仿射 Lie 代数对称性} % 原文: 4.2 Quantization: Virasoro And Affine Lie Algebra Symmetry

我们已经看到，群流形上的粒子具有 $G_L\times G_R$ 对称性。在空间取为 $S^1$ 的经典二维理论中，在耦合常数的特殊值 $f^2=12\pi c_1k$ 处，它被提升为无穷维的 $(LG)_L\times(LG)_R$ 对称性。

本节我们要证明：圆上的\emph{量子}二维 WZW 理论具有 $\widetilde{LG}_L\times\widetilde{LG}_R$ 对称性，其中 $\widetilde{LG}$ 是环路群（loop group）$LG=Map(S^1,G)$ 的一个中心扩张。

$LG$ 的 Lie 代数是 $Lie(LG)=L\mathfrak{g}$，即 $S^1$ 映入 $\mathfrak{g}$ 的映射构成的向量空间，它有一个 $\mathbb{R}$ 的中心扩张。我们将把 Lie 代数中的环路按一组基 $T_a$ 分解为：
\begin{equation}
  J^a(x) = \sum_{n\in\mathbb{Z}}J^a_n e^{2\pi i n x}
  \tag{4.51}
  \label{eq:4.51}
\end{equation}
其中我们选取 Lie 代数的一组正交归一（ON）基 $T_a$。现在我们取空间为一个圆周，且 $x\sim x+1$。

我们的约定是：$su(N)$ 是\emph{反厄米}（antihermitian）矩阵构成的 Lie 代数，$so(N)$ 是\emph{反对称}矩阵构成的 Lie 代数，等等。然而，我们希望保持标准的厄米性关系 $(J^a_n)^\dagger=+J^a_{-n}$。因此我们的流将取值于厄米矩阵的空间，等等。

定义 $\widetilde{L\mathfrak{g}}$ 的结构常数是：$k$ 是中心的，且
\begin{equation}
  [J^a_n, J^b_m] = i f^{abc}J^c_{n+m} + k\delta^{ab}n\,\delta_{n+m,0}
  \tag{4.52}
  \label{eq:4.52}
\end{equation}
这被称为一个\emph{仿射 Lie 代数}（affine Lie algebra）。

请特别注意，取 $n=0$ 时，生成元 $J^a_0$ 张成一个同构于 $\mathfrak{g}$ 的有限维子代数。

我们将勾勒出在 WZW 模型中推导这一对称性的三种途径。为确定起见，取 $S=kS^+_{WZW}$。

\textbf{第一种途径：}

按 PW 公式，$g\to ge^{\epsilon}$ 下的 Noether 流为
\begin{equation}
  J^R = J_- = i(12\pi c_1 k)g^{-1}\partial_- g
  \tag{4.53}
  \label{eq:4.53}
\end{equation}
类似地，对 $g\to e^{-\epsilon}g$ 有
\begin{equation}
  J^L = J_+ = -i(12\pi c_1 k)\partial_+ g\,g^{-1}
  \tag{4.54}
  \label{eq:4.54}
\end{equation}
其中因子 $i$ 来自上文解释过的我们的约定。
% [Tutor Note: (4.53)(4.54) 与 (4.44)(4.45) 相差整体常数：这里 $J^R=J_-$、$J^L=J_+$ 指的是同一个守恒流的两种归一化，$i(12\pi c_1k)$ 这个因子只是把 (4.32) 中 $kS^+_{WZW}$ 的整体系数与「几何」定义的 $g^{-1}\partial_- g$ 联系起来，不影响守恒律；但它决定了流代数里能级 $k$ 出现在哪里，以及 (4.59) 的 Schwinger 项系数。另一处易错点是 $i$ 的来源：按本段约定 $su(N)$ 取反厄米矩阵，而要求 $(J^a_n)^\dagger=+J^a_{-n}$，故 Hermitian 流与反 Hermitian 生成元之间必须插入一个 $i$；(4.53) 与 (4.54) 的 $i$ 符号相反，正对应 $\partial_-g$ 与 $\partial_+g$ 的左右之别，切勿写成同号。]

由于运动方程 $\partial_+J^R_-=0$ 与 $\partial_-J^L_+=0$，实际上存在一个无穷维对称性
\begin{equation}
  \delta^R_\epsilon : g\to ge^{\epsilon(x^-)}
  \tag{4.55}
  \label{eq:4.55}
\end{equation}
\begin{equation}
  \delta^L_\epsilon : g\to e^{\epsilon(x^+)}g
  \tag{4.56}
  \label{eq:4.56}
\end{equation}

我们在运动方程的层面上已经研究过这一点。在量子理论中，这些变换由下式生成
\begin{equation}
  \delta^R_\epsilon\mathcal{O} = \left[\int dx^-\,\epsilon^a(x^-)J^a_-(x^-), \mathcal{O}\right]
  \tag{4.57}
  \label{eq:4.57}
\end{equation}
其中 $\mathcal{O}$ 是任意算符。左对称性有类似的表达式。

流本身的变换很容易算出：
\begin{equation}
  \begin{aligned}
    \delta^R_\epsilon J^R &= -\left[\epsilon(x^-), J^R(x^-)\right] + i(12\pi c_1 k)\epsilon'(x^-)\\
    \delta^R_\epsilon J^L &= 0
  \end{aligned}
  \tag{4.58}
  \label{eq:4.58}
\end{equation}

把它与 (4.57) 结合，我们推出：
\begin{equation}
  [J^a(x^-), J^b(y^-)] = i f^{abc}J^c(y^-)\delta(x^--y^-) + i(12\pi c_1 k)\delta^{ab}\frac{d}{dx^-}\delta(x^--y^-)
  \tag{4.59}
  \label{eq:4.59}
\end{equation}

当我们按 (4.51) 定义流并使用 $c_1=1/(24\pi^2)$ 时，我们看到 (4.59) 恰好等价于 (4.52)。
% [Tutor Note: 这步「恰好等价」是本节最容易翻车的地方，务必自己推一遍。取 $c_1=1/(24\pi^2)$ 得 $12\pi c_1k=k/(2\pi)$，故 (4.59) 的正则项为 $i\frac{k}{2\pi}\delta^{ab}\partial_{x^-}\delta(x^--y^-)$；作逆变换 $J^a_n=\int_0^1dx\,J^a(x)e^{-2\pi i n x}$（与 (4.51) 的 $e^{+2\pi i nx}$ 配套），这个 $\delta'$ 项给出的能级项正比于 $n\,\delta_{n+m,0}$，系数取决于把 $\partial_{x^-}\delta(x^--y^-)$ 的导数算在 $x^-$ 上还是算在 $y^-$ 上（二者差一个符号），也取决于 (4.51) 的 Fourier 指数取 $+2\pi i nx$ 还是 $-2\pi i nx$。正因为如此，(4.59) 与后面 (4.72) 的 $\delta'$ 项看起来差一个整体符号，而两处描述的其实是同一个流代数。核对时请以「能级 $k$ 的符号与大小」为准绳，并参照 Witten 的论文固定约定。]

\textbf{第二种途径：}

第二种量子化的途径是使用光锥形式。让我们把 $x^+$ 视为时间、$x^-$ 视为空间。于是，在光锥形式中，作用量 $S^+_{WZW}$ 对「时间」（即 $x^+$）的导数是一阶的。因此它已经处于哈密顿形式。现在我们再一次运用关于「对时间导数为一阶的作用量」的那条一般性评述。

变分作用量得到：
\begin{equation}
  \delta\left(kS^+_{WZW}\right) = -12\pi c_1 k\int dx^+\left[\int dx^-\operatorname{Tr}\left(g^{-1}\delta g\right)\partial_+\left(g^{-1}\partial_- g\right)\right]
  \tag{4.60}
  \label{eq:4.60}
\end{equation}

与 (??) 比较，我们得到辛形式
\begin{equation}
  \Omega = \int dx^-\operatorname{Tr}\left(g^{-1}\delta g\right)\partial_-\left(g^{-1}\delta g\right) - g^{-1}\partial_- g\left(g^{-1}\delta g\right)^2
  \tag{4.61}
  \label{eq:4.61}
\end{equation}
% [原文标记，原文 p.111：此处交叉引用在源文件中未解析，印作 (??)；据上下文应指前文「把一阶作用量写成 $\int p\,dq-H\,dt$ 并由此读出辛形式」的那组一般性公式]
% [Tutor Note: (4.61) 按印本转写，但印本里 $\operatorname{Tr}$ 后面没有括号，看起来只管第一项；从上下文（以及 (4.65) 中 $L+\xi g^{-1}dg$ 与 $\xi=-6\pi c_1k$ 的对应）判断，迹应当覆盖两项：$\Omega=\int dx^-\operatorname{Tr}\left[(g^{-1}\delta g)\partial_-(g^{-1}\delta g)-g^{-1}\partial_-g(g^{-1}\delta g)^2\right]$。另外 (4.60)→(4.61) 省掉的步骤是：由 $\delta(kS^+)=-\int dx^+dx^-\operatorname{Tr}\left[(g^{-1}\delta g)\partial_+(g^{-1}\partial_-g)\right]\cdot 12\pi c_1k$ 作一次分部积分（把 $\partial_+$ 移到 $g^{-1}\delta g$ 上），再读出 $\Omega=\delta\Theta$ 型的一次形式；这也正是后面「把 (4.61) 写成 $\delta$ 下的全导数」那道练习要补的一步。]

由此可以计算流的 Poisson 括号。这正是 Witten 的论文中所做的，读者可参阅该文了解细节。

{\small\itshape ［原文标记：关于光锥哈密顿量量子化之缺点的评注：你会遗漏沿 $x^+=\mathrm{const}$ 行进的波。在这个形式里它们就是「空间切片」，但仍然是左行自由度，因此在另一个时间切片中具有动力学重要性。］}
% [原文标记，原文 p.112 页首，通篇以上下两行星号包裹；内容为完整的评注而非待办，故全文译出后保留为小号斜体]

\medskip
\noindent\textbf{练习}\quad 把 (4.61) 写成 $\delta$ 下的全导数。

\textbf{第三种途径：}

这里最好回顾一下关于群流形 $G$ 上粒子的哈密顿形式的那番讨论。

我们可以不取光锥规范而直接进入哈密顿形式，但这有一点微妙。设 $\delta$ 表示在固定时刻的场构形空间 $LG=Map(S^1,G)$ 上的微分。

$LG$ 在环路 $g(\theta)$ 处的切空间同构于 $L\mathfrak{g}$：
\begin{equation}
  T_g LG \cong L\mathfrak{g}
  \tag{4.62}
  \label{eq:4.62}
\end{equation}
{\small\itshape ［原文标记：在 $S^1$ 上沿空间方向的坐标究竟记作 $x$ 还是 $\theta$，要前后一致］}
% [原文待办标记，原文 p.112 页边，以 \ensuremath{\clubsuit} 号包裹]

$L\mathfrak{g}$ 的对偶空间是 $\Omega^1(S^1;\mathfrak{g}^*)$，因此作为 $T^*LG$ 纤维上一个坐标的典范动量 $L$ 属于
\begin{equation}
  L \in \Omega^1(S^1;\mathfrak{g}^*)
  \tag{4.63}
  \label{eq:4.63}
\end{equation}
其配对为
\begin{equation}
  \langle L, \epsilon\rangle = \oint d\theta\,\langle L(\theta), \epsilon(\theta)\rangle
  \tag{4.64}
  \label{eq:4.64}
\end{equation}

我们用不变形式把 $\mathfrak{g}\cong\mathfrak{g}^*$ 等同起来，从而可以把 $T^*LG$ 上的辛形式写成
\begin{equation}
  \Omega = \oint_{S^1}\operatorname{Tr}\delta L\,g^{-1}\delta g + \operatorname{Tr}\left(L + \xi g^{-1}dg\right)\left(g^{-1}\delta g\right)^2
  \tag{4.65}
  \label{eq:4.65}
\end{equation}
其中 $\xi=-6\pi c_1 k$。这里 $g=g(\theta)$ 是 $G$ 中的环路，而 $g^{-1}\delta g$ 是一个取值于 $L\mathfrak{g}^*\cong L\mathfrak{g}$ 的 $LG$ 上的 1-形式，并且是 $S^1$ 上的 0-形式。$L$ 现在被视作取值于 $L\mathfrak{g}$ 的 $S^1$ 上的一个 1-形式。因此 (4.65) 是 $T^*LG$ 上的一个 2-形式。证明它在 $\delta$ 下封闭是一个很好的练习。

现在，同样地，要写出哈密顿形式的 Legendre 变换，我们需要找到 $T^*G$ 上的一个 1-形式，使得 $\Omega=\delta\Theta$。为此我们需要把 $g$ 延拓为圆盘的一个映射 $\tilde g\in Map(D^2,G)$。于是可以写：
\begin{equation}
  \Theta = \oint_{S^1}\operatorname{Tr}\left(Lg^{-1}\delta g\right) - \xi\int_{D^2}\operatorname{Tr}\left(\tilde g^{-1}d\tilde g\right)^2\left(\tilde g^{-1}\delta\tilde g\right)
  \tag{4.66}
  \label{eq:4.66}
\end{equation}

那么可以验证
\begin{equation}
  S = \int_\gamma\Theta - \int dt\,6\pi c_1 k\left[\oint_{S^1}\frac{1}{2}\operatorname{Tr}L^2 + \frac{1}{2}\operatorname{Tr}\left(g^{-1}\partial_x g\right)^2\right]
  \tag{4.67}
  \label{eq:4.67}
\end{equation}
确实给出 WZW 作用量。

确实，这正是我们在圆盘上量子化 Chern-Simons 理论所得到的形式！见上文 (3.86)。

让我们计算 Poisson 括号。这番讨论紧密追随群流形上点粒子的情形。见上文方程 (??)。按 Lie 代数的一组正交归一基写出，我们有
\begin{equation}
  \Omega = \frac{1}{2}\oint dx\oint dy
  \begin{pmatrix}\delta L_a(x) & \delta\epsilon^a(x)\end{pmatrix}
  \begin{pmatrix}
    0 & \delta^a{}_b\,\delta(x-y)\\[4pt]
    -\delta_a{}^b\,\delta(x-y) & -f^c{}_{ab}\left(L_c + \xi g^{-1}\partial_y g\right)\delta(x-y)
  \end{pmatrix}
  \begin{pmatrix}\delta L_b(y)\\ \delta\epsilon^b(y)\end{pmatrix}
  \tag{4.68}
  \label{eq:4.68}
\end{equation}
% [原文标记，原文 p.113：此处交叉引用在源文件中未解析，印作 (??)；据上下文应指前文群流形上点粒子的 Poisson 括号那组公式（即 $T^*G$ 上的 $\Omega^{AB}$）]

于是逆矩阵为：
\begin{equation}
  \Omega^{AB} =
  \begin{pmatrix}
    -f^c{}_{ab}\left(L_c + \xi\left(g^{-1}\partial_y g\right)_c\right)\delta(x-y) & -\delta_a{}^b\,\delta(x-y)\\[4pt]
    \delta^a{}_b(x-y) & 0
  \end{pmatrix}
  \tag{4.69}
  \label{eq:4.69}
\end{equation}

由 (4.69) 得到的 Poisson 括号是
\begin{equation}
  \begin{aligned}
    \{g(x), g(y)\} &= 0\\
    \{L_a(x), g(y)\} &= -g(y)T_a\,\delta(x-y)\\
    \{L_a(x), L_b(y)\} &= -f^c{}_{ab}\left(L_c(y) + \xi\left(g^{-1}g\right)_c(y)\right)\delta(x-y)
  \end{aligned}
  \tag{4.70}
  \label{eq:4.70}
\end{equation}
% [Tutor Note: (4.70) 第三行印作 $(g^{-1}g)_c(y)$，对照 (4.68)(4.69) 中同一位置应为 $(g^{-1}\partial_y g)_c(y)$——否则该项是常数、与 $y$ 的依赖和 (4.62) 中「切空间」的几何意义都不成立，故几乎肯定是排印时漏掉了 $\partial_y$；此处按印本原样转写，不作改动。另需注意 (4.68)(4.69) 中指标的位置：$\delta^a{}_b$ 与 $\delta_a{}^b$ 分别出现在右上、左下两块，二者正是互为转置的「单位张量」，若把上下指标写平（都写成 $\delta^a_b$）就会看不出矩阵的反对称结构，从而在求逆时丢掉一个符号。]

对 $L_a$ 变分作用量，在时间演化下给出 $L_a=6\pi c_1 k(g^{-1}\dot g)$，从而
\begin{equation}
  J^a(x) = L^a(x) - 6\pi c_1 k\left(g^{-1}\partial_x g\right)^a
  \tag{4.71}
  \label{eq:4.71}
\end{equation}
是一个右行流。

现在，利用 (4.70) 我们可以容易地算出 $J^a(x)$ 的 Poisson 括号：
\begin{equation}
  \{J^a(x), J^b(y)\} = -f^{abc}J^c(y)\delta(x-y) - \frac{k}{2\pi}\delta^{ab}\partial_x\delta(x-y)
  \tag{4.72}
  \label{eq:4.72}
\end{equation}
于是恢复了流代数。

\textbf{Remark（评注）：} 利用与左对称性共轭的角动量，我们可以类似地导出左行流的流代数。

\medskip
\noindent\textbf{练习}\quad 证明确实有 $\Omega=\delta\Theta$。

提示：解释为什么
\begin{equation}
  \delta\operatorname{Tr}\left(g^{-1}dg\right)^2\left(g^{-1}\delta g\right) + d\operatorname{Tr}\left(g^{-1}\delta g\right)^2\left(g^{-1}dg\right) = 0
  \tag{4.73}
  \label{eq:4.73}
\end{equation}
